\documentclass[border=10pt]{standalone}
\usepackage{fontspec}
\setmainfont{Noto Serif}
\newfontfamily\bengalifont{Noto Serif Bengali}[Script=Bengali]
\usepackage[Latin,Bengali]{ucharclasses}
\setTransitionsFor{Bengali}{\begingroup\bengalifont}{\endgroup}
\usepackage{newunicodechar}
\newunicodechar{।}{{\bengalifont\char"0964}}
\newunicodechar{॥}{{\bengalifont\char"0965}}

\usepackage[table]{xcolor}
\usepackage{booktabs,tabularx,array}
\definecolor{bapilrule}{HTML}{1F2937}
\definecolor{bapilstripe}{HTML}{F3F4F6}
\renewcommand{\arraystretch}{1.28}
\setlength{\tabcolsep}{5pt}

\begin{document}
\begin{minipage}{17cm}
{\large\bfseries সারণি ৬১।}\ {\large এই পর্যালোচনার আগের খসড়ার সঙ্গে নির্দেশক বাজেটের মিলকরণ (কোটি টাকা, ৩৬ মাস)}\par\vspace{7pt}
\small
\rowcolors{2}{bapilstripe}{white}
\begin{tabularx}{\linewidth}{@{}>{\hsize=1.974\hsize\raggedright\arraybackslash}X>{\hsize=0.513\hsize\raggedright\arraybackslash}X>{\hsize=0.513\hsize\raggedright\arraybackslash}X@{}}
\arrayrulecolor{bapilrule}\toprule
\rowcolor{white}\textbf{ধাপ} & \textbf{৩৬ মাসের মোট} & \textbf{পরিবর্তন} \\
\midrule
পর্যালোচনার আগের খসড়া বাজেট (অক্টোবর ২০২৬) & ১৮.৪৯ & সূচনাবিন্দু \\
দ্বিতীয় ও তৃতীয় বছরের অঙ্ক প্রথম বছরের মূল্যে পুনর্নির্ধারণ & ১৮.২৯ & -০.২০ \\
টাকায় পরিশোধ্য ব্যয়ে সুস্পষ্ট মূল্যবৃদ্ধি (বছরে ৮\%) & ১৯.৭৯ & +১.৫০ \\
প্রথম বছরের ধাপে ধাপে নিয়োগ (গড়ে ৭৫\% পদ পূরণ) & ১৯.২৯ & -০.৫০ \\
মূল্যায়নের নতুন উপাদান (আওতাবহির্ভূত প্রশ্নসেট, উদ্ধৃতি-বিচার, ব্যবস্থাগুলোর বার্ষিক নিরীক্ষা) & ১৯.৫৪ & +০.২৫ \\
\textbf{প্রত্যক্ষ ব্যয়ের ৫\% পৃথক আপৎকালীন সংস্থান = সংশোধিত মূল হিসাব} & \textbf{২০.৪৩} & +০.৮৯ \\
\bottomrule
\end{tabularx}
\par\vspace{5pt}{\footnotesize\itshape প্রতিটি ধাপ blueprint/\allowbreak{}src/\allowbreak{}budget\_model.py (bridge ফাংশন) দ্বারা একবারে একটি অনুমান বদলে হিসাব করা হয়েছে। খসড়ার বার্ষিক মোট (৪.৯৬, ৬.৭১ ও ৬.৮৩) যোগ করলে ১৮.৫০ হতো কেবল প্রতিটি আলাদাভাবে আসন্নীকৃত হওয়ায়; আসন্নীকরণের আগে খসড়ার মোট ছিল ১৮.৪৮৬৮ কোটি টাকা।}
\par\vspace{3pt}{\scriptsize\color{gray} উৎস: BAPIL কারিগরি ব্লুপ্রিন্ট, সংস্করণ ১.০ (২০২৬)। CC BY 4.0।}
\end{minipage}
\end{document}
